\documentclass[11pt]{ctexart}
\usepackage[a4paper,top=2cm,bottom=2cm,left=2.2cm,right=2.2cm]{geometry}
\usepackage{fancyvrb}
\usepackage{xcolor}
\usepackage{colortbl}
\usepackage{fontspec}
\usepackage[most]{tcolorbox}
\usepackage{booktabs}
\usepackage{enumitem}
\usepackage{titlesec}
\usepackage{hyperref}

\setmonofont{Consolas}
\setCJKmonofont{SimHei}
\setlength{\parindent}{0pt}
\setlist{nosep,leftmargin=1.4em}

% ── Step3 主题色（与 Step1/2 同族、换为深青） ──
\definecolor{deep}{RGB}{14,84,88}        % 主色：深青
\definecolor{accent}{RGB}{28,128,120}    % 强调：青绿
\definecolor{keyline}{RGB}{90,130,200}   % 知识点：蓝
\definecolor{keybg}{RGB}{238,244,255}
\definecolor{warnline}{RGB}{200,120,60}  % 警告：橙
\definecolor{warnbg}{RGB}{255,243,235}
\definecolor{okline}{RGB}{26,116,88}     % 成果：绿
\definecolor{okbg}{RGB}{240,252,244}
\definecolor{tipline}{RGB}{88,158,148}   % 提示：青
\definecolor{tipbg}{RGB}{238,248,246}
\definecolor{codebg}{RGB}{249,249,252}   % 代码底
\definecolor{rowtint}{RGB}{236,246,244}  % 表头浅青

\titleformat{\section}{\large\bfseries\color{deep}}{\thesection}{0.6em}{}
\titleformat{\subsection}{\normalsize\bfseries\color{accent}}{\thesubsection}{0.5em}{}
\titleformat{\part}[display]{\centering\Huge\bfseries\color{deep}}{}{0pt}{\Huge}
\titlespacing{\section}{0pt}{14pt}{6pt}
\titlespacing{\subsection}{0pt}{10pt}{4pt}

\newtcolorbox{keybox}[1]{enhanced, colback=keybg, colframe=keyline,
  boxrule=0.6pt, arc=1.5mm, left=3mm, right=3mm, top=2mm, bottom=2mm,
  title={\bfseries #1}, fonttitle=\small\bfseries, coltitle=white,
  colbacktitle=keyline}
\newtcolorbox{warnbox}[1]{enhanced, colback=warnbg, colframe=warnline,
  boxrule=0.6pt, arc=1.5mm, left=3mm, right=3mm, top=2mm, bottom=2mm,
  title={\bfseries #1}, fonttitle=\small\bfseries, coltitle=white,
  colbacktitle=warnline}
\newtcolorbox{okbox}[1]{enhanced, colback=okbg, colframe=okline,
  boxrule=0.6pt, arc=1.5mm, left=3mm, right=3mm, top=2mm, bottom=2mm,
  title={\bfseries #1}, fonttitle=\small\bfseries, coltitle=white,
  colbacktitle=okline}
\newtcolorbox{tipbox}[1]{enhanced, colback=tipbg, colframe=tipline,
  boxrule=0.5pt, arc=1.5mm, left=3mm, right=3mm, top=1.5mm, bottom=1.5mm,
  title={\bfseries #1}, fonttitle=\small\bfseries, coltitle=white,
  colbacktitle=tipline}

\newcommand{\code}[1]{\texttt{#1}}
\DefineVerbatimEnvironment{Verilog}{Verbatim}{frame=single,fontsize=\small}

\begin{document}

% ==================== 封面 ====================
\begin{titlepage}
\centering
\vspace*{2.2cm}
{\color{deep}\Huge\bfseries 云台闭环控制}\\[6pt]
{\color{deep}\Large\bfseries 从零手搓 Step3 教程}\\[14pt]
{\large 2026 FPGA 创新设计赛道（安路）赛题二}\\[6pt]
{\large 视觉方向 · Step3\_closed\_loop}\\[20pt]

\begin{tcolorbox}[enhanced, colback=codebg, colframe=accent, boxrule=0.8pt,
  arc=2mm, width=12.5cm, left=4mm, right=4mm, top=3mm, bottom=3mm]
\centering
{\small 读者：零基础队员（建议先读完 Step1、Step2 教程）\\
目标：让云台自己把目标"追"回画面中心——从坐标到运动\\
前置要求：会 Verilog 基础语法、跑过 ModelSim、懂 Step2 的输出\\
预计学习时间：10～14 天（每天 2～3 小时）}
\end{tcolorbox}

\vspace{8mm}
\begin{tcolorbox}[enhanced, colback=tipbg, colframe=tipline, boxrule=0.5pt,
  arc=2mm, width=12.5cm, left=4mm, right=4mm, top=2mm, bottom=2mm]
\centering
{\small\color{deep} \textbf{本教程与你手上的代码是什么关系}\\
Step2 的输出是"目标坐标"；Step3 把它变成云台的\textbf{转动}——\\
误差计算、比例控制、死区、限速限幅、舵机 PWM，一套闭环。}
\end{tcolorbox}

\vfill
{\small\color{black!50} fpga-ptz-tracker 教程系列 · 2026-09}
\end{titlepage}

\tableofcontents
\newpage

% ============================================================
\part{从零学习教程}

% ============================================================
\section{你在造什么}

\subsection{一句话}

\textbf{把"目标在画面哪里"变成"云台往哪转多少"，让镜头自动追着目标走。}

Step2 每帧告诉你：目标中心在 \code{(cx, cy)}。Step3 要回答：舵机脉宽该加
多少微秒——左偏 100 像素，脉宽该从 1500 变成多少？

\subsection{闭环长什么样}

\begin{Verbatim}
        +----------+     误差(px)    +----------+   脉宽(us)   +----------+
画面 -->| 检测链   | -------------> | 控制律   | ----------> | PWM 舵机 |
        | (Step2)  |                | (本步)   |             | (本步)   |
        +----------+                +----------+             +----+-----+
             ^                                                   |
             |            云台带着相机转，目标回到中心             |
             +---------------------------------------------------+
\end{Verbatim}

这是一个\textbf{反馈闭环}：误差决定动作，动作改变误差。这一节教你写
"大脑"（\code{gimbal\_ctrl}，控制律）和"肌肉"（\code{pwm\_servo}，舵机驱动）。

\subsection{做完是什么效果}

上电后云台在中位静止 → 屏幕角落看到红色目标偏左 → 云台\textbf{平滑地}
转过去，几秒内目标稳稳停在画面中心小方框里，不再抖动。
配合 Step4 的 OSD，画面上还能显示当前状态（LOCK / LOST、坐标、误差）。

\begin{okbox}{学完 Step3 你能做到}
\begin{itemize}
  \item 说清闭环的四个要素：测量、误差、控制、执行
  \item 手算 Q8.8 定点比例控制的一帧步进
  \item 解释死区、限速、限幅各防什么事故
  \item 用增益公式判断一组 Kp 会不会振荡
  \item 跑通 3 个测试（步进表、PWM 测量、闭环在环收敛）
\end{itemize}
\end{okbox}

% ============================================================
\section{阶段 0：反馈控制入门（1～2 天）}

\subsection{开环 vs 闭环}

\begin{keybox}{两种控制方式}
\textbf{开环}：只发命令不看结果。像蒙着眼走路——步子迈出去就收不回来，
地面稍有起伏就摔。\\
\textbf{闭环}：先看结果和目标的差距（误差），再决定下一步动作。
误差大就大步走，误差小就小步走，到了就停。
\end{keybox}

云台跟踪天生就是闭环：目标会动、云台有机械误差、装配有偏差，
这些都不用精确知道——只要每帧看一眼"偏了多少"，闭环自己会修正。

\subsection{一个生活类比：调淋浴水温}

你想把水温调到 40 度（目标），手拧阀门（执行），皮肤测温（测量）：

\begin{itemize}
  \item 水太凉 → 往热拧一点；太烫 → 往冷拧一点（比例控制）
  \item 快到 40 度了 → 只敢轻轻拧（误差小，步子小）
  \item 在 39～41 度之间 → 不动了（死区，防止来回抖）
  \item 阀门拧到底了 → 不再硬拧（限幅，行程保护）
\end{itemize}

\begin{warnbox}{新手最常见的直觉错误}
"误差大就猛打"——听起来合理，实际会\textbf{振荡}：你猛拧到 40 度，
但热水管道有延迟，过两秒温度冲到 45 度；你又猛拧回来，又冲到 35 度……
永远停不下来。\textbf{闭环里的"延迟"决定了你敢用多大的步子}，
这是本教程最重要的一个知识点（阶段 3 展开）。
\end{warnbox}

\subsection{练习}

\begin{enumerate}
  \item 用"调水温"的类比，说出测量、误差、控制、执行分别对应什么
  \item 想一个开环控制会失败、闭环能救回来的场景
  \item 思考：为什么"步子大"在无延迟时很好、有延迟时会振荡？
\end{enumerate}

% ============================================================
\section{阶段 1：舵机与 PWM（2 天）}

\subsection{舵机怎么听话}

本项目用 SG90/MG90S 级别的位置舵机：三根线（电源、地、信号），
信号是一路 50Hz 的 PWM 脉冲，\textbf{脉冲宽度}决定转角：

\begin{Verbatim}
0.5ms ----- 0 度
1.0ms ----- 45 度
1.5ms ----- 90 度（中位）
2.0ms ----- 135 度
2.5ms ----- 180 度
周期固定 20ms（50Hz），每个周期只发一个脉冲
\end{Verbatim}

换算成比例：180 度 / 1000us = \textbf{0.18 度/us}。这个数后面算增益要用，
先记住。

\subsection{PWM 怎么生成：计数器 + 比较}

FPGA 生成 PWM 只需要一个计数器和一个比较器：

\begin{Verilog}
// 50MHz 时钟：20ms = 1,000,000 拍；1us = 50 拍
period_cnt <= (period_cnt >= PERIOD-1) ? 0 : period_cnt + 1;
pwm = (period_cnt < duty_us * 50);   // 前 duty 微秒输出高电平
\end{Verilog}

想输出 1500us 脉冲？让 \code{pwm} 在计数器前 75000 拍为高。
分辨率是 1 拍 = 20ns = 0.02us，比舵机的机械精度高两个数量级，绰绰有余。

\begin{keybox}{为什么"共周期"是硬要求}
水平、垂直两个舵机\textbf{必须共用同一个周期计数器}。如果各用各的，
两个计数器相位慢慢漂移，X 轴脉宽落点和 Y 轴错开，供电电流互相拉扯，
舵机抖动。\code{pwm\_servo} 里一个计数器出两路比较，结构性杜绝。
\end{keybox}

\subsection{上电安全：默认中位}

复位后 duty 寄存器就是 1500——云台上电即指向正前方，不会猛甩。
这是比赛硬约束"冷启动直接跟踪"的一部分：固化进 bit 流的参数
让板子一上电就处于可演示状态。

\subsection{练习}

\begin{enumerate}
  \item 手算：想让云台转到 120 度，脉宽多少 us？对应多少拍（50MHz）？
  \item 手算：1500us 脉宽 = 多少拍？
  \item 思考：如果时钟是 100MHz，1us 等于多少拍？分辨率提升多少？
\end{enumerate}

% ============================================================
\section{阶段 2：比例控制与定点算术（2～3 天）}

\subsection{控制律：一帧只做五步}

每个"帧拍"（每秒 60 次）执行一遍，以水平轴为例：

\begin{Verilog}
1) err  = cx - 640                    // 误差：目标偏离中心多少像素
2) s    = (err * KP) >>> 8            // 比例：像素 -> us，Q8.8
3) 死区：|err| <= 25  ->  s = 0        // 目标在中心附近不动
4) 限速：|s| > 100    ->  s = ±100     // 单帧最多加 100us
5) 限幅：duty+s 超出 700~2300 -> 截断   // 行程硬保护
最终：duty = duty_old + s
\end{Verilog}

五个步骤的顺序不能乱：先死区（免算）、再限速（防猛）、后限幅（防越界）。

\subsection{Q8.8：用移位代替小数}

KP 想表达"每像素动 0.078 微秒"这种小数，FPGA 不做浮点。
约定：参数里存的是\textbf{真值 × 256}（高 8 位整数、低 8 位小数）：

\begin{Verilog}
KP = 20                    // 参数值（寄存器里存的数）
真值 = 20 / 256 ≈ 0.078 us/px
s = (err * KP) >> 8        // 先乘大数，再右移 8 位除掉 256
\end{Verilog}

\begin{keybox}{为什么右移必须用 >>>（算术右移）}
err 是\textbf{有符号}数（目标可以在中心左边）。Python 的 \code{>>}
对负数是向下取整，Verilog 的 \code{>>>} 配合 \code{signed} 才是同一语义。
用 \code{>>} 或漏 \code{signed}，目标在左半边时行为就错——
Step2 里 rgb2hsv 已经吃过一次这个亏，这里原样复用纪律。
\end{keybox}

\subsection{手算一遍（本阶段练习核心）}

取 KP=20、死区 25、限速 100、限幅 700～2300、当前 duty=1500：

\begin{Verilog}
例 A：cx = 840      err = 200
      s = (200*20)>>8 = 4000>>8 = 15   duty: 1500 -> 1515
例 B：cx = 660      err = 20   （死区内）
      s = 0                            duty 不动
例 C：cx = 1240     err = 600
      s = (600*20)>>8 = 46             duty: 1500 -> 1546
例 D：cx 在极远处   err = 20000（假设）
      s 本应 1562，限速截到 100        duty: 1500 -> 1600
\end{Verilog}

\begin{warnbox}{中间量位宽}
\code{err(13bit 有符号) × KP(16bit)} 的积最多 29 bit——
task 里中间量显式声明 \code{signed [28:0]}。位宽不够会静默截断，
误差大的时候步进值直接变成垃圾。
\end{warnbox}

\subsection{练习}

\begin{enumerate}
  \item 手算：err = -150（目标在左），KP=20，duty=1500，新的 duty？
  \item 手算：err = 500 时若 KP=90，s 是多少？和 KP=20 差几倍？
  \item 用 Python 验证你的答案（注意负数 >> 的行为）
\end{enumerate}

% ============================================================
\section{阶段 3：延迟与稳定性——为什么 Kp 不能大（3 天）}

这一阶段是 Step3 的灵魂。控制律本身 20 行就写完，
\textbf{参数选错，云台当场表演抽搐}——我们在闭环仿真里真的复现了一次
（跑飞到 109 度），这里教你提前算出来。

\subsection{延迟从哪来}

从"目标偏了"到"画面反映云台已转"，要经过：

\begin{Verilog}
帧 N:   检测链算出 err（1 帧流水）           -> 帧末锁存坐标
帧 N:   控制律算出 step，写 duty              -> 执行
帧 N+1: 舵机转到新角度，画面已变              -> 采集
帧 N+1: 检测链才看到"转完之后"的画面          -> 下一次误差
\end{Verilog}

合计约 \textbf{2 帧延迟}（1 帧检测 + 1 帧执行）。延迟期间你是在
"闭着眼睛"继续输出——这是振荡的根源。

\subsection{手推一次振荡（不用任何公式）}

假设 Kp 很大，看误差怎么一步步演化（每帧 2 帧延迟，比例全量响应）：

\begin{Verilog}
帧:   0     1     2     3     4     5     6     7
err:  200   200   -200  -200  +200  +200  -200  -200
      ^^^^^^^^^^^ 误差永远在 ±200 之间来回跳 —— 等幅振荡！
\end{Verilog}

Kp 小一点呢？每次只修正误差的一半：200 → 100 → 50 → 25 → 死区，收敛。
\textbf{Kp 大到什么程度会开始振荡？}下面用增益算。

\subsection{增益公式（现场换硬件后要重算的唯一样子）}

三个系数连乘：

$$L = \underbrace{KP/256}_{\text{us/px}} \times \underbrace{0.18}_{\text{度/us}} \times \underbrace{20.6}_{\text{px/度}} \approx KP \times 0.0144 \quad (\text{度/度/帧})$$

再乘延迟放大（2 帧延迟约 ×2），得到\textbf{有效增益}，稳定判据：

\begin{Verilog}
2 × L < 1    即    KP < 256 / (2 × 0.18 × 20.6) ≈ 34
\end{Verilog}

代入我们整定时的三组参数验证：

\begin{center}
\begin{tabular}{llll}
\toprule
\rowcolor{rowtint}\textbf{KP} & \textbf{单帧增益 L} & \textbf{有效增益 2L} & \textbf{实测结果} \\
\midrule
90   & 1.30 & 2.60 & 发散：角度单调跑飞到 109 度 \\
40   & 0.58 & 1.16 & 慢发散振荡（>1，越振越大） \\
\textbf{20}   & \textbf{0.29} & \textbf{0.58} & \textbf{12 帧锁定，无极限环} \\
\bottomrule
\end{tabular}
\end{center}

\textit{修订注记（2026-10-03）：上表三组实验是初版结构的整定记录（回路按 2 帧延迟）。
后续控制节拍修正（tick 改为"新目标包到达拍"，即 T1）把回路延迟降到 1 帧，
判据相应放宽为 $L<1$；\textbf{默认 KP=20 不变，实测 settle\_frame 由 36 改善到 12}。
下面的 ×2 推导保留原样，作为"延迟决定可用增益"的完整示范。}

\begin{keybox}{为什么是"×2"而不是"×1"或"×3"}
2 帧延迟的离散比例反馈，每周期误差翻转一次并放大 L 倍，
两个帧才完成一个往返——所以是 ×2。若延迟是 1 帧（比如云台响应极快），
判据放宽到 L<1；延迟 3 帧则收紧到 L<1/3。
\textbf{延迟越大，可用增益越小}——这也解释了为什么"检测链要快"。
\end{keybox}

\begin{warnbox}{限速救不了错误的 Kp}
有人想"Kp 大了加限速就行"。限速能压住单帧冲击，但压不住
符号交替的振荡（每次都 ±100us 来回摆）。先保证 2L<1，再谈限速。
\end{warnbox}

\subsection{练习}

\begin{enumerate}
  \item 手算：KP=30 时的 L 和 2L，判断稳不稳
  \item 若换一个 FOV 90 度的相机（px/度 = 1280/90 约 14.2），KP 上限变多少？
  \item 思考：检测链流水加深 1 帧，KP 上限怎么变？
\end{enumerate}

% ============================================================
\section{阶段 4：安全设计与状态字（1～2 天）}

\subsection{四道保险}

\begin{center}
\begin{tabular}{lll}
\toprule
\rowcolor{rowtint}\textbf{保险} & \textbf{参数} & \textbf{防什么} \\
\midrule
死区   & \code{DZ=25px}      & 目标在中心附近时云台抖动 \\
限速   & \code{SLEW=100us/帧} & 单帧猛打，机械冲击 \\
软限位 & \code{700～2300us}  & 云台撞机械极限（控制律内） \\
硬限位 & \code{500～2500us}  & 软限位失效后的第二道（PWM 内） \\
\bottomrule
\end{tabular}
\end{center}

\subsection{丢失目标怎么办}

某帧检测失败（\code{target\_valid=0}）：\textbf{保持当前脉宽不动}，
只把 \code{lost\_cnt} 加一。目标可能只是被遮挡了一瞬，
猛回中或猛保持旧方向都可能错——保持是最安全的选择。
连续丢失很多帧的应对策略（搜索/回中）留给上层，底层只报告。

\subsection{状态字 O\_status}

8 位状态字，每一位都给 OSD 或调试用：

\begin{Verilog}
O_status[0] = dz_x    水平在死区内（目标已在中心）
O_status[1] = dz_y    垂直在死区内
O_status[2] = lost    本帧丢失目标
O_status[3] = lim_x   水平限幅命中（打到边界了）
O_status[4] = lim_y   垂直限幅命中
\end{Verilog}

\code{track\_en} 拨码断开时：强制回中位（1500），用于演示复位/安全停车。

\subsection{练习}

\begin{enumerate}
  \item 为什么要"双重限幅"？只有 PWM 里那道行不行？
  \item 画出"丢失目标 3 帧后重新发现目标"时 duty 的时间线
  \item 思考：lost 时保持脉宽，什么情况下这个策略会出问题？
\end{enumerate}

% ============================================================
\section{阶段 5：跑通闭环测试（1～2 天）}

\subsection{一键回归}

\begin{Verbatim}
cd D:\fpga-ptz-tracker\tools
python run_all.py tb_gimbal tb_pwm tb_closed_loop
\end{Verbatim}

三个测试各管一层，从里到外：

\begin{center}
\begin{tabular}{lll}
\toprule
\rowcolor{rowtint}\textbf{测试} & \textbf{验证什么} & \textbf{规模} \\
\midrule
tb\_gimbal     & 控制律 10 步进表（死区/限速/限幅/丢失） & 10 帧 \\
tb\_pwm        & PWM 波形逐拍测量 & 2 个完整周期 \\
tb\_closed\_loop & 相机-舵机模型在环，60 帧内锁定 & 全闭环 \\
\bottomrule
\end{tabular}
\end{center}

\subsection{怎么读结果}

\begin{Verbatim}
--- tb_gimbal ...
    # STEP3 OK duty_x=05E6 duty_y=05D7 lost=0000 dz_x=0 dz_y=1
    ==> PASS          ← 10 步全对（含死区、限速、限幅、丢失保持）

--- tb_pwm ...
    # PERIOD=1000000 pulse_x=75000 pulse_y=75000   ← 上电中位
    ==> PASS          ← 周期与脉宽逐拍精确

--- tb_closed_loop ...
    TB_RESULT: PASS settle_frame=12 final_err=-8
    ==> PASS          ← 12 帧锁定（±25px 死区内保持 10 帧）
\end{Verbatim}

\begin{okbox}{Step3 学习清单（全部打勾就过关）}
\begin{itemize}
  \item 能说清闭环四要素，以及"每帧一小步"为什么够用
  \item 能手算一帧步进（含死区/限速/限幅三种截断）
  \item 能解释 Q8.8 和 \code{>>>} 的配合
  \item 能用增益公式判断一组 Kp 的稳定性，说出 KP 上限怎么来的
  \item 能说出四道保险各防什么、丢失目标的策略
  \item 跑通 3 个测试并读懂每一行输出
\end{itemize}
\end{okbox}

% ============================================================
% Step3 教程 · Part II：工程实现详解
% 由 step3_tutorial.tex \input 引入（共用同一份导言区与主题色）
% ============================================================

\part{工程实现详解}

\begin{tipbox}{这一部分写什么}
Part I 讲了知识；这一部分把仓库里\textbf{真实做出来的 Step3 工程}逐模块
讲透：每个参数为什么是这个值、每个赋值为什么用阻塞/非阻塞、
三轮 Kp 整定是怎么在闭环仿真里收敛的。对应目录：
\code{03\_vision/step3\_closed\_loop/}——2 个 RTL（209 行）、2 个 TB、
3 轮整定记录。
\end{tipbox}

% ============================================================
\section{工程总览}

\subsection{文件清单}

\begin{center}
\begin{tabular}{llr l}
\toprule
\rowcolor{rowtint}\textbf{文件} & \textbf{类型} & \textbf{行数} & \textbf{职责} \\
\midrule
\code{gimbal\_ctrl.v} & RTL & 140 & 控制律：死区+比例+限速+限幅，双轴 \\
\code{pwm\_servo.v} & RTL & 69 & 双路 50Hz PWM，共周期，上电中位 \\
\code{tb/tb\_gimbal.v} & 测试 & 62 & 控制律 10 步进表逐字段比对 \\
\code{tb/tb\_pwm.v} & 测试 & 94 & PWM 波形逐拍连续监视器（v2） \\
\code{sim/gimbal\_stim.txt} & 向量 & 10 & 帧输入序列 \\
\code{sim/gimbal\_exp.txt} & 向量 & 10 & 金标准期望 \\
\bottomrule
\end{tabular}
\end{center}

\subsection{数据流与时钟域}

\begin{Verilog}
Step2 检测链 (74.25MHz 像素域)
   |  目标坐标 + frame_tick (每帧 1 拍)
   v
pulse_sync  (单 bit 脉冲跨域, Step1 模块)
   v
gimbal_ctrl (50MHz PWM 域)  err/step/duty 计算
   |  O_duty_x/y_us (13bit, us)
   v
pwm_servo   (50MHz)         双路 50Hz 波形 -> 舵机
\end{Verilog}

两个时钟域：像素域 74.25MHz 跑检测，PWM 域 50MHz 跑控制与波形。
跨域只传两样东西：\textbf{帧拍脉冲}（pulse\_sync，单 bit）和
\textbf{目标坐标}（cdc\_4phase，多 bit）。50MHz 是刻意选择的：
1us 恰好 50 拍，µs 到拍的换算纯乘法无余数。

\subsection{两条设计原则（本步新增）}

\begin{keybox}{Step3 的方法论}
\textbf{1. 无符号纪律}：误差、乘积、累加一律显式 \code{\$signed}；
task 内中间量 29bit。整步最阴的 bug 全部出在符号上。\\
\textbf{2. 参数即固化}：所有控制参数是 RTL parameter——综合进 bit 流，
上电即生效，满足"冷启动直接跟踪"的比赛硬约束。
\end{keybox}

% ============================================================
\section{模块详解：gimbal\_ctrl}

\subsection{task 复用：双轴一份逻辑}

水平、垂直两轴的控制律完全同构，只有参数不同。RTL 用一个 task
\code{calc\_axis} 承载全部算术，时序块里调两次：

\begin{Verilog}
task calc_axis(input signed [12:0] err, input [15:0] kp, ...);
    prod = err * $signed({1'b0, kp});   // 29bit 有符号积
    s    = prod >>> 8;                  // Q8.8 -> us, floor
    ...
endtask
\end{Verilog}

注意 \code{\$signed(\{1'b0, kp\})}：位拼接默认产生无符号数，
不包一层 \code{\$signed} 的话整个乘法退化成无符号，负误差直接回绕——
这正是闭环发散事故的模型侧根因（详见振荡复盘一节）。

\subsection{五步顺序与实现}

\begin{Verilog}
dz_hit = (err >= -dz) && (err <= dz);      // 1) 死区判定
if (dz_hit) s = 0;                          //    命中则步进清零
if (s >  SLEW_P)  s =  SLEW_P;              // 2) 限速（双向）
else if (s < -SLEW_P) s = -SLEW_P;
nx = $signed({1'b0, duty_in}) + s;          // 3) 累加
if (nx > SMAX_P) begin duty_out = SMAX_P; lim_hit = 1; end
else if (nx < SMIN_P) begin duty_out = SMIN_P; lim_hit = 1; end
\end{Verilog}

\subsection{输出同拍锁存"新值"（一个真 bug）}

初版输出寄存写成 \code{O\_duty\_x\_us <= duty\_x}——而 \code{duty\_x}
在同一拍才被赋成新值，非阻塞赋值读的是\textbf{旧值}，输出滞后一帧。
tb\_gimbal 的步进表第一帧就对不上。修复：

\begin{Verilog}
O_duty_x_us <= (I_track_en && I_target_valid) ? dnx_x : duty_x;
//                              ^^^^^ 本拍刚算出的新值
\end{Verilog}

\begin{warnbox}{非阻塞赋值的"读旧写新"}
非阻塞赋值 RHS 立即求值、LHS 帧末更新。同一个 always 块里
"先更新 \code{duty\_x}，再输出 \code{duty\_x}"，读到的是更新前的值。
要么像本工程一样输出 \code{dnx\_x}（新值），要么想清楚就是要旧值。
\end{warnbox}

\subsection{状态字与丢失策略}

\begin{Verilog}
O_status <= {2'b00, limy, limx, lost, dzy, dzx};   // 一拍拼装
// lost 帧: duty 保持不动, O_lost_cnt 饱和累加
// track_en=0: duty 强制回中位, 状态清零
\end{Verilog}

一个迭代 1 抓出的细节：无效帧分支里 \code{dzx/limx} 清零用非阻塞，
同拍组装的 \code{O\_status} 采样到\textbf{旧值}——状态字残留上一帧的死区位。
修复：该分支的清零改\textbf{阻塞赋值}（ 组装在同一个顺序块里，
先清后拼）。这是"阻塞/非阻塞选择不是教条、而是看数据依赖"的实例。

% ============================================================
\section{模块详解：pwm\_servo}

\subsection{µs 到拍的换算}

\begin{Verilog}
localparam TICKS_PER_US = 50;               // 50MHz / 1e6
wire [19:0] pulse_x = duty_x * TICKS_PER_US; // 1500us -> 75000 拍
\end{Verilog}

20bit 计数器数到 1,000,000（50Hz 周期）回零，两路输出
\code{pwm = (period\_cnt < pulse)}——\textbf{一个计数器、两个比较器}，
共周期是结构保证的。

\subsection{第二道限幅与上电中位}

\begin{Verilog}
wire [12:0] dx = (I_duty_x_us > 13'd2500) ? 13'd2500 :
                 (I_duty_x_us < 13'd500)  ? 13'd500  : I_duty_x_us;
// 复位值: duty_x <= 13'd1500    —— 上电即中位
\end{Verilog}

控制律已经限在 700～2300，这里再限 500～2500——两道保险的区间
故意不同：如果上层逻辑跑飞给出荒谬值，第二道还能兜住机械极限。

% ============================================================
\section{测试平台设计}

\subsection{tb\_gimbal：10 步进表覆盖全部分支}

金标准 \code{gen\_tb\_gimbal()} 构造 10 帧输入序列，
\textbf{刻意覆盖每种截断}：

\begin{center}
\begin{tabular}{lll}
\toprule
\rowcolor{rowtint}\textbf{帧} & \textbf{输入} & \textbf{考察点} \\
\midrule
0-1 & cx=840 (err=200) & 常规比例步进（1500 到 1515） \\
2   & cx=900 (err=260) & 累加继续 \\
3   & cx=700 (err=60)  & 小误差小步进 \\
4   & cx=641 (err=1)   & \textbf{死区命中，duty 不动} \\
5-6 & valid=0          & \textbf{丢失保持 + lost\_cnt 累加} \\
7   & cx=100 (err=-540) & 负方向 + 大步进 \\
8   & cx=640 (err=0)   & 死区 \\
9   & cx=660 (err=20)  & 死区 \\
\bottomrule
\end{tabular}
\end{center}

期望文件每行 \code{\{dx:03X\}\{dy:03X\}\{lost:04X\}\{dzx:X\}\{dzy:X\}}，
TB 逐字段比对——duty、lost 计数、两个死区标志\textbf{都要对}。
第一条期望 \code{5EB5DC000001} 手工可验：dx=0x5EB=1515（err=200 步进 15）、
dy=0x5DC=1500（cy=360 正中死区）、dz\_y=1。

\subsection{tb\_pwm：从任务式同步到连续监视器}

第一版 TB 用任务式 \code{@(posedge tick)} 等待周期起点再测量——
实测 \code{len=1}：TB 的等待事件和 DUT 的置位\textbf{同拍竞态}，
采样点漂移。v2 重写为\textbf{连续监视器}：

\begin{Verilog}
// 自由运行计数器 + tick 锁存，没有任何 @ 事件同步
always @(posedge clk) begin
    if (pwm) pulse_len <= pulse_len + 1;
    if (tick) begin
        report(frame_cnt, pulse_len, period_len);
        pulse_len <= 0;
    end
    period_len <= (tick) ? 1 : period_len + 1;
end
\end{Verilog}

锁存窗口\textbf{包含 tick 拍本身}（v1.1 曾差一拍测得 999999），
换挡激励避开过渡周期。v2 测得：上电中位 75000/75000/1000000 逐拍精确，
换挡后 60000（1200us）/100000（2000us）精确。

\subsection{tb\_closed\_loop：相机-舵机模型在环}

真正的验收是把整条链闭环跑起来。TB 内建行为级相机模型
（参数都有实物依据）：

\begin{center}
\begin{tabular}{lll}
\toprule
\rowcolor{rowtint}\textbf{参数} & \textbf{取值} & \textbf{依据} \\
\midrule
光学放大 & 20 px/度 & 62.2 度 FOV @ 1280px \\
舵机增益 & 0.18 度/us & 180 度 / 1000us 行程 \\
机械速率限幅 & 1.5 度/帧 & 约 90 度/s \\
系统延迟 & 1 帧 & 检测链流水；控制节拍 = 新目标包到达拍（T1 修正前为 2 帧） \\
初始指向 theta0 & 30 度 & 与目标差 10 度 \\
\bottomrule
\end{tabular}
\end{center}

模型每个 PWM 周期算 \code{theta\_cmd = 30 + (duty-1500)×0.18}，
带速率限幅逼近；TPG 方块按 \code{x = 640 + (40-theta)×20} 实时移动
——\textbf{相机真的在"看"云台转完的画面}。验收标准：
60 帧内误差进入 ±25px 死区并保持 10 帧。

% ============================================================
\section{验证结果与振荡事故复盘}

\subsection{回归矩阵}

\begin{center}
\begin{tabular}{lll}
\toprule
\rowcolor{rowtint}\textbf{测试} & \textbf{比对粒度} & \textbf{结果} \\
\midrule
tb\_gimbal & 10 步进表逐字段（duty/lost/dz 标志） & PASS \\
tb\_pwm & 波形逐拍（脉宽/周期） & PASS \\
tb\_closed\_loop & 在环收敛（settle\_frame） & PASS：12 帧锁定 \\
\bottomrule
\end{tabular}
\end{center}

\subsection{Step3 抓到的真缺陷}

\begin{center}
\begin{tabular}{llp{5.4cm}l}
\toprule
\rowcolor{rowtint}\textbf{编号} & \textbf{现象} & \textbf{根因} & \textbf{类型} \\
\midrule
G1 & 输出滞后一帧 & 输出寄存同拍读旧值（非阻塞语义） & RTL \\
G2 & 状态字残留 & 无效帧清零用非阻塞，同拍拼装读旧值 & RTL \\
G3 & 测得脉宽 len=1 & 任务式 @ 同步与 DUT 置位同拍竞态 & TB \\
G4 & 周期 999999 & 计数窗口不含 tick 拍 & TB \\
G5 & 闭环跑飞 109 度 & 模型无符号减法回绕 + 回路增益超标 & 模型+参数 \\
\bottomrule
\end{tabular}
\end{center}

\subsection{振荡事故复盘（为什么值得单独一节）}

整定第一轮 Kp=90/70，闭环仿真 60 帧里角度\textbf{单调跑飞到 109 度}，
duty 打满限幅。复盘出两个叠加的问题：

\begin{warnbox}{事故 1：TB 模型的无符号减法}
相机模型算误差时用了无符号减法：\code{(40 - theta)} 在 theta 超过 40
时回绕成巨大正数，模型认为目标突然飞到另一边——闭环被喂了假误差。
修复：全链路 \code{\$signed}。教训：\textbf{行为模型也要遵守无符号纪律}，
它错的时候 RTL 怎么对都没用。\\[4pt]
\textbf{事故 2：回路增益超标}。Kp=90 时单帧增益 L=90×0.0144=1.30，
两帧延迟放大后 2L=2.60——每往返误差翻 2.6 倍，单调发散。
第二轮 Kp=40（2L=1.16）仍是慢发散；第三轮 Kp=20/18（2L=0.58）进入稳定收敛。
\textbf{闭环专项修正（T1）}：控制节拍改为"新目标包到达拍"，回路延迟由 2 帧降到
1 帧（判据放宽为 L<1），实测 \textbf{settle\_frame 由 36 改善到 12}，
RTL 默认参数同步更新。
\end{warnbox}

\begin{keybox}{给现场整定者的操作清单（完整版见 IF-SERVO-PWM.md）}
1）先验极性：给固定小误差，云台必须\textbf{朝减小误差}的方向动；
反了改 \code{INVERT\_X/Y} 参数，不要倒接电机线。\\
2）从保守 Kp 起步（按增益公式算出上限，取一半）。\\
3）逐步加大直到出现轻微振铃，回退 30\%。\\
4）SLEW 从 100 起步：只影响"最多多快"，不影响稳定性判据。\\
5）改完参数重新综合固化——参数是 RTL parameter，不烧不生效。
\end{keybox}

% ============================================================
\section{上板挂载说明（给上板组）}

\begin{center}
\begin{tabular}{ll}
\toprule
\rowcolor{rowtint}\textbf{项目} & \textbf{内容} \\
\midrule
时钟 & 50MHz PWM 域（官方例程已有该域，直接挂） \\
跨域输入 & 目标坐标（cdc\_4phase）；控制节拍本地生成 = 新目标包到达拍
（T1 后不再用 pulse\_sync 传 tick） \\
输出 & O\_pwm\_x/y 接舵机信号线；O\_duty\_dbg 供 ILA 观测 \\
按键/拨码 & track\_en 接拨码（0 = 回中位，演示/安全用） \\
供电 & 舵机独立 5V 供电，\textbf{与 FPGA 共地不共电} \\
参数 & KP/DZ/SLEW/限幅全部是 parameter，现场整定后重编译固化 \\
\bottomrule
\end{tabular}
\end{center}

上板三步：
\begin{enumerate}
  \item 两个 RTL 加入 TD 工程，实例化 gimbal\_ctrl + pwm\_servo；
  \item 把 Step2 的 \code{O\_frame\_tick}/坐标经跨域模块接入；
  \item 先不接舵机，用 ILA 看 O\_pwm 波形确认周期/脉宽，再带载。
\end{enumerate}

\begin{tipbox}{和 Step2 的分工}
Step2 只负责"看"，输出坐标和帧拍；Step3 只负责"动"，输入坐标和帧拍。
两者通过两个跨域信号解耦——检测链改坏了不影响舵机安全（丢失保持），
舵机故障也不阻塞检测。\textbf{接口即边界}，约定见 \code{06\_interfaces}。
\end{tipbox}

% ============================================================
\section{小结：Step3 的可复用产出}

\begin{itemize}
  \item \textbf{2 个可综合 RTL 模块}（209 行），控制律与波形生成解耦；
  \item \textbf{一套增益设计方法}：L = KP×0.0144，按回路延迟选判据
    （初版 2 帧延迟用 2L<1；T1 修正后延迟 1 帧，判据放宽为 L<1），
    换相机/舵机后重算即可；
  \item \textbf{一个在环验证范式}：行为级相机模型 + TPG 方块跟随，
    无硬件就能整定参数；
  \item \textbf{三份整定记录}：发散/慢发散/锁定的完整过程，
    是答辩时"工程严谨性"的最好素材。
\end{itemize}

\begin{okbox}{下一步（Step4 预告）}
坐标有了、云台动起来了，但屏幕上还看不到任何"系统在干活"的痕迹。
Step4 把结果画回画面：状态文字、目标框、跨域目标包、整链联调——
评委会看到的画面，从这里来。
\end{okbox}


% ============================================================
\part*{结语}
\clearpage
\section*{总结与下一步}
\addcontentsline{toc}{section}{总结与下一步}

\subsection*{Step3 学到了什么}
\addcontentsline{toc}{subsection}{Step3 学到了什么}
\begin{itemize}
  \item 反馈闭环的思想：测量、误差、控制、执行
  \item 舵机 PWM 规格（50Hz、0.5～2.5ms、0.18 度/us）与共周期生成
  \item 控制律五步：误差、Q8.8 比例、死区、限速、限幅
  \item 稳定性判据：增益公式与 2 帧延迟的 ×2 放大
  \item 安全设计：四道保险、丢失保持、状态字、上电中位
  \item 参数整定方法论：先极性、后增益、在环仿真验证
\end{itemize}

\subsection*{下一步学什么}
\addcontentsline{toc}{subsection}{下一步学什么}
\begin{itemize}
  \item \textbf{Step4 OSD 与系统联调}：把检测结果画回画面
  \begin{itemize}
    \item 字库 ROM 与文字叠加（状态行直接显示 LOCK/坐标）
    \item 68bit 目标包跨域（视觉域 → 显示域）
    \item vision\_top 整链集成与全帧逐位比对
  \end{itemize}
  \item 对应文档：\code{03\_vision/step4\_osd\_integration/}
  \item 对应金标准：\code{tools/golden\_model.py::gimbal\_ctrl\_step}
\end{itemize}

\end{document}
